\section{Background and related work}
\label{sec:background}

\subsection{Digital Twins in manufacturing}

Digital Twin research in manufacturing has developed around virtual representations of physical assets, physical--digital connectivity, lifecycle data, simulation, monitoring, and synchronisation \cite{grieves2017digital,kritzinger2018digital,tao2018digital,fuller2020digital}. Reviews show that manufacturing remains a central Digital Twin domain, with recurring attention to production systems, predictive maintenance, optimisation, cyber-physical integration, and interoperability \cite{negri2017review,semeraro2021digital}. These works provide the technical baseline for the present paper: a Digital Twin is not merely a static model, but a connected digital representation that can support understanding, prediction, and decision-making.

However, the literature is less mature on the front-end problem of how the knowledge needed for a useful Digital Twin is elicited from work practice. Manufacturing Digital Twin papers often focus on asset data, simulation models, control architectures, or integration levels. Human knowledge is frequently assumed to enter through requirements, domain expertise, dashboard design, or validation, but the transformation from situated operator explanation to structured Digital Twin knowledge is rarely the main object of study.

\subsection{Human-in-the-loop and human-centred Digital Twins}

Human-in-the-loop Digital Twin research argues that Digital Twins should not be understood as fully autonomous replacements for human expertise. Instead, they should distribute observation, interpretation, decision-making, and responsibility between digital systems and human actors. Work on human roles in Digital Twins highlights that decisions about automation, human oversight, trust, and responsibility are often under-specified \cite{agrawal2023humans,parasuraman2000levels,lee2004trust}. Human-centred smart-factory research similarly emphasises that shop-floor workers remain central to knowledge creation, interpretation, and operational adaptation \cite{longo2022ubiquitous}.

For Notebook-to-OSL, this literature justifies treating operators as knowledge-bearing participants in Digital Twin development rather than only as end users. It also shows why the output of elicitation must preserve responsibility, escalation, and validation needs. A recommendation or alarm in a Digital Twin is only useful if the system also represents when an operator should trust it, when it should be checked, and who is responsible for acting.

\subsection{Operator knowledge, situation awareness, and expert judgement}

The knowledge used by operators is often partly tacit. Polanyi's account of tacit knowledge is commonly summarised by the idea that people can know more than they can explicitly tell \cite{polanyi1966tacit}. In manufacturing, this is visible when operators recognise abnormal behaviour through experience-based cues that are difficult to specify in advance. Such cues may include sound, vibration, smell, timing, tool marks, alarm combinations, or the relation between machine response and operation context.

Operator decision-making also depends on situation awareness and control context. Endsley describes situation awareness in terms of perceiving elements in the environment, comprehending their meaning, and projecting future status \cite{endsley1995situation}. Rasmussen's distinction between skill-, rule-, and knowledge-based behaviour shows that action in technical systems cannot be reduced to isolated commands \cite{rasmussen1983skills}. Together, these perspectives imply that a Digital Twin should not only capture that an operator acts, but also what the operator noticed, how the cue was interpreted, and why the action was judged appropriate.

\subsection{Situated learning, cognitive apprenticeship, and contextual inquiry}

Situated learning theory argues that learning is embedded in social and material practice rather than detached from context \cite{lave1991situated}. Cognitive apprenticeship makes a related argument in design terms: expert reasoning can be made visible through modelling, coaching, scaffolding, articulation, reflection, and exploration \cite{collins1989cognitive}. This is directly relevant to the proposed training situation. When an operator teaches a researcher how a CNC process is understood and handled, the operator is not only transmitting information; they are demonstrating how expertise is enacted in context.

Contextual inquiry provides a more operational method base. Its classic orientation is to study work where it happens, in partnership with the practitioner, and to treat interpretations as things that should be checked rather than silently inferred \cite{beyer1995apprenticing,beyer1997contextual}. Contextual design also shows that field observations can be transformed into structured work models. Notebook-to-OSL adopts this logic but narrows the target: rather than modelling the whole work system or immediately designing a user interface, it captures operator-strategy fragments that can later inform Digital Twin artefacts.

\subsection{Cognitive task analysis and hidden decision knowledge}

Contextual inquiry is strong on work setting, artefacts, flow, and interpretation, but it is less specific about hidden cognitive demands. Cognitive task analysis (CTA) and the critical decision method focus more directly on expertise, decision points, cues, anomalies, difficult judgments, and expert--novice differences \cite{crandall2006workingminds,klein1989critical}. ACTA-based work has also shown how CTA can be customised for design and requirements-oriented projects, where expert cognition must be elicited and converted into design-relevant outputs \cite{stanik2021acta}.

Notebook-to-OSL uses CTA as a clarification layer. During live training, the researcher should avoid turning the situation into a formal interview. After the session, however, selected notes can be probed with CTA-inspired questions: What first indicated that this was abnormal? What would a novice miss? What else could this cue mean? When would the action be wrong? When should the situation be escalated? These questions help recover cognitive structure without overloading the live training session.

\subsection{Requirements engineering, MBSE, and model-ready transformation}

Requirements engineering and systems engineering add the rigor needed to turn field material into engineering artefacts. Requirements work distinguishes stakeholder statements, needs, rationale, risks, constraints, verification criteria, and formal specifications \cite{vanlamsweerde2009requirements}. SEBoK guidance similarly emphasises stakeholder identification, source traceability, rationale, risks, status, ownership, and transformation from stakeholder needs to system requirements \cite{sebokStakeholder,sebokRequirements}. These principles justify why the Notebook-to-OSL schema includes fields such as source, rationale, uncertainty, risk, validation need, responsibility, and status.

The downstream modelling literature is also relevant. EARS provides lightweight controlled natural-language patterns for transforming messy requirements statements into more consistent requirement-like structures \cite{mavin2009ears}. SysML v2 provides a formal MBSE target for requirements, structure, behaviour, analysis, and verification artefacts \cite{omgSysml2024}. Semantic Digital Twin work using ontologies and knowledge graphs shows that Digital Twin knowledge is increasingly represented as reusable and interoperable structures rather than isolated documents \cite{karabulut2023ontologies,listl2024knowledge,braun2022graph}. These works motivate the term \emph{model-ready}: the fragment does not need to be fully formal, but it must contain enough structure to support later mapping to requirements, OSL, SysML v2, rules, or semantic models.

\subsection{Synthesis and research gap}

The literature provides strong components but not the complete method needed here. Contextual inquiry and situated learning support learning from real work in context. CTA supports probing cues and expert decisions. Requirements engineering and MBSE support traceability, rationale, risk, and transformation. Semantic Digital Twin literature supports model-ready knowledge structures. What remains underdeveloped is the bridge between these elements: a practical workflow for converting situated operator-training notes into traceable, uncertainty-aware, validated, model-ready operator-strategy fragments for human-in-the-loop Digital Twin development. Notebook-to-OSL is proposed as that bridge.
